


%%%%%%%%%%%%%%%%%%%%%%%%%%%%%%%%%%%%%%%%%%%%%%%%%%%%%%%%%%%%%%
%%%%%%%%%%%%% MA-XXXX Introducción a las Demostraciones %%%%%%%%%%%%%
%%%%%%%%%%%%%%%%%%%%%%%%%%%%%%%%%%%%%%%%%%%%%%%%%%%%%%%%%%%%%%
\newpage
\subsection*{MA-02XX Introducción a las Demostraciones}
\addcontentsline{toc}{subsection}{MA-02XX Introducción a las Demostraciones}

\begin{refsection}
\begin{table}[H]
\centering{
\begin{tabular}{|ll|}
\hline
%\multicolumn{2}{|c|}{\textbf{Nivel de virtualidad del curso:} Virtual}\\ \hline
\textbf{Año:} I  & \textbf{Requisitos:} MA-1XXX Matemática Exploratoria \\ 
\textbf{Ciclo:} II  & \textbf{Correquisitos:} Ninguno \\ 
\textbf{Tipo de curso:} Teórico & \textbf{Horas presenciales por semana:}   \\ 
\textbf{Créditos:} 4 & \textbf{Horas de trabajo independiente por semana:}  \\ \hline
\end{tabular}
}
\end{table}


\subsubsection*{Descripción del curso}

Este curso se ubica en el II ciclo y está dirigido al estudiantado de la carrera de Matemática que ha aprendido las destrezas del curso MA-XXX1 Matemática Introductoria I.

Se describe este curso como la segunda parte de un bloque introductorio de dos cursos, que aporta a la formación del estudiantado el nivel de matemática superior necesario para abordar los desafíos de los cursos posteriores, al tiempo que el curso tiene el reto de  despertar curiosidad, motivación y pasión por la matemática. El estudiantado conocerá un proceso donde la curiosidad lo enfrentará con inquietudes particulares que se responderán a lo largo de su carrera.


\subsubsection*{Objetivos}

El objetivo principal del curso es continuar desarrollando las habilidades operacionales del estudiante a través de temas matemáticos básicos y necesarios para el quehacer cotidiano.
Continuar desarrollando las habilidades operacionales del estudiante a través de contenidos matemáticos básicos, requeridos en cursos posteriores.[vRAZR]
Extender el desarrollo de habilidades operacionales del estudiantado mediante contenidos matemáticos básicos, requeridos en cursos posteriores. [vRAZR usando Bloom]

\begin{enumerate}
\item Reproducir procedimientos en los contenidos matemáticos del curso que luego permitan al estudiantado interpretar, descubrir, inferir y hasta predecir argumentaciones y conexiones dentro de la matemática.

\item Estimular (fomentar) en el estudiantado capacidades de análisis e interpretación que permitan el desarrollo de habilidades operacionales y el primer estímulo de pensamiento abstracto, ambos conceptos piedras angulares presentes en los mecanismos de razonamiento posteriores.
\end{enumerate}

\subsubsection*{Contenidos}

\begin{enumerate}
\item \textbf{Contenido 1:} 

\item \textbf{Contenido 2:} 

\item \textbf{Contenido 3:} 

\item \textbf{Contenido 4:} 

\item \textbf{Contenido 5:} 
\end{enumerate}

\subsubsection*{Metodología}



\printbibliography[heading=subbibliography,section=\the\value{refsection}]
%\begin{bibunit}
%\putbib[Biblio]
%\end{bibunit}

\end{refsection}
